% GROK REPORT::CONTRACT1::2026-10-04T12:56:00-07:00
% !TEX program = xelatex
\documentclass[11pt]{article}
\usepackage[margin=0.88in]{geometry}
\usepackage{amsmath,amssymb}
\usepackage{hyperref}

\title{GROK note --- standard SAT self-reduction does\\not contract a polynomial exponent\\[0.25em]
\large CONTRACT1}
\author{Agent P$\neq$NP / GROK explorer}
\date{2026-10-04 12:56 PT}

\begin{document}
\maketitle

P versus NP remains open, and this note does not prove a time lower bound for 3-SAT.

\section*{Status}
\textbf{OPEN.} The candidate squeeze was: a uniform charged map taking any exact polynomial 3-SAT algorithm of unknown finite exponent $K$ to another exact algorithm of exponent at most $a+bK$ with $0<b<1$, repeatable, so every finite $K$ is sent to the attractor $a/(1-b)$; if that fixed point lay in an unconditional forbidden region of the \emph{same} model, that would be a contradiction. No such $a,b$ are known. This note tries the missing step and records the obstruction. It does not claim $P\neq NP$. Root \texttt{main.tex} was not edited. No GitHub Actions were added.

Three requirements, each marked as it stands after the attempt:
\begin{enumerate}
\item \textbf{Uniform charged contraction with $0<b<1$.} \textbf{OPEN.} Not found. The standard downward self-reduction has $b=1$ (\S\ref{sec:map}). Treating that self-reduction as a contraction is \textbf{FALSIFIED/RETIRED}.
\item \textbf{Repeatable, so the attractor is $a/(1-b)$.} The fixed-point algebra is \textbf{PROVED/EXACT} as a recurrence (\S\ref{sec:algebra}), not as a SAT bound. As a step in a 3-SAT contradiction it is \textbf{OPEN}, blocked on the missing map.
\item \textbf{Unconditional forbidden neighborhood in the same model.} \textbf{OPEN.} Not available. Williams and Fortnow--Lipton--van Melkebeek--Viglas restrict space. Fortnow--Santhanam and Dell--van Melkebeek rule out strong shrinkage only \textbf{CONDITIONAL} on a polynomial-hierarchy collapse, and they rule out compression, not every non-shrinking transformation (\S\ref{sec:lit}).
\end{enumerate}
Status of the contraction: \textbf{OPEN}. Status of P versus NP: \textbf{OPEN}.

\section*{What was opened}
Primary or standard secondary pages actually fetched for this note:
\begin{itemize}
\item Fortnow--Santhanam, ECCC TR07-096 (abstract).
\item Dell--van Melkebeek, STOC 2010 paper (author PDF, including Theorem~1 and Corollary~1).
\item Impagliazzo--Paturi--Zane, FOCS 1998 / JCSS 2001 author PDF (Sparsification Lemma, Corollary~1).
\item Fortnow--van Melkebeek, ECCC TR00-028 (abstract); Fortnow--Lipton--van Melkebeek--Viglas, JACM 2005 abstract on van Melkebeek's page.
\item Williams, ECCC TR07-036 (abstract); Williams, ``Some Open Problems Regarding Lower Bounds for NP'' (Gasarch open-problems column).
\item Williams, STOC 2024, ``Self-Improvement for Circuit-Analysis Problems'' (author PDF), checked because it is the nearest self-improvement theorem and does \emph{not} supply the missing map.
\end{itemize}
Searches for a uniform charged map $K\mapsto a+bK$ with $0<b<1$ for exact 3-SAT decision, all costs charged, did not return one. No citation is invented.

\section{Exponent map of standard self-reduction}
\label{sec:map}

\subsection*{Setup}
Let $m$ be the bit length of a SAT instance and let $n\le m$ be its number of variables. Let $A$ be an exact decision algorithm with
\[
T_A(m)\le C\, m^{K}
\]
for some constant $C\ge 1$ and some finite exponent $K\ge 1$. Restricting one variable does not increase bit length: each restricted instance still has length at most $m$.

\subsection*{The $n$-call reduction sends $K$ to $K+1$}
\textbf{PROVED/EXACT} (algebra of this recurrence only; not a SAT lower bound).

The usual search reduction fixes variables one at a time. For each of the $n$ variables, query $A$ on the formula with that variable set to $0$; if the answer is satisfiable keep $0$, otherwise keep $1$. That is at most $n\le m$ calls, each of cost at most $C m^{K}$, plus $O(m)$ work to write each restriction. The charged cost is
\begin{align*}
T(m)
&\le n\cdot C m^{K} + n\cdot O(m) \\
&\le C\, m\cdot m^{K} + O(m^{2}) \\
&= O\bigl(m^{K+1}\bigr),
\end{align*}
using $K\ge 1$ so that the $O(m^{2})$ term is absorbed. The exponent map is therefore
\[
K\mapsto K+1 = 1 + 1\cdot K,
\]
that is $a=1$ and $b=1$ in the shape $a+bK$. In particular $b\not< 1$.

One iteration step does not contract. Repeating the same charging gives
\[
K_{t+1}=K_{t}+1,\qquad K_{t}=K_{0}+t.
\]
As $t\to\infty$, $K_{t}\to\infty$. The attractor formula $a/(1-b)$ is undefined at $b=1$. This is not a contraction.

The two-call decision split (satisfiable if and only if at least one of the two one-variable restrictions is) charges
\[
T(m)\le 2\, C m^{K}+O(m)=O(m^{K})
\]
if the polynomial algorithm is invoked once on each branch and not unfolded. The degree is unchanged. Unfolding the recurrence $T(n)\le 2T(n-1)+p(n)$ down to constant size produces an exponential, which is a worse upper bound, not a smaller polynomial degree.

\subsection*{Fixed-point algebra, if a contraction existed}
\label{sec:algebra}
\textbf{PROVED/EXACT} as algebra, not as a property of SAT. If a sequence obeyed $K_{t+1}\le a+b K_{t}$ with constants $a$ and $0<b<1$, then
\[
K_{t}\le a\frac{1-b^{t}}{1-b}+b^{t} K_{0}
\xrightarrow[t\to\infty]{}\frac{a}{1-b}.
\]
No $a,b$ making this a charged uniform map on exact polynomial 3-SAT are known, so the limit is not a 3-SAT exponent.

\textbf{EXAMPLE only, not a SAT exponent.} The toy parameters $a=1$, $b=1/2$, $K_{0}=6$ were iterated exactly (Python \texttt{fractions}). The values are $6$, $4$, $3$, $5/2$, $9/4$, $17/8$, $33/16$, $65/32$, and the limit is $1/(1-1/2)=2$. The same code with $b=1$, $a=1$, $K_{0}=6$ produced $6,7,8,9,10,11$ over five steps and diverges. These numbers are \textbf{COMPUTATIONAL} checks of the recurrence. They are not running times of a SAT algorithm.

\section{A failed attempt at a contraction}
\label{sec:fail}

\textbf{PROVED/EXACT} (algebra). Suppose one hoped that a single call on $n-1$ variables would force a smaller degree, via
\[
T(n)\le n^{c}+T(n-1)
\]
with $c\ge 0$ fixed. It does not.

For integers $n\ge 1$ and $K\ge 1$,
\[
(n-1)^{K}=n^{K}\Bigl(1-\frac{1}{n}\Bigr)^{K}\le n^{K}\Bigl(1-\frac{1}{n}\Bigr)=n^{K}-n^{K-1},
\]
because $0\le 1-1/n\le 1$ and the exponent $K\ge 1$ does not increase the value. Hence if $C\ge 1$ and $K\ge c+1$,
\begin{align*}
n^{c}+C(n-1)^{K}
&\le n^{c}+C n^{K}-C n^{K-1} \\
&\le C n^{K},
\end{align*}
since $C n^{K-1}\ge n^{c}$ whenever $c-(K-1)\le 0$ and $C\ge 1$. A polynomial upper bound of degree $K\ge c+1$ is compatible with the recurrence. The recurrence does not force a degree $bK$ with $b<1$.

Solving the recurrence by unfolding gives only the upper bound $T(n)\le T(0)+\sum_{i=1}^{n} i^{c}\le T(0)+n\cdot n^{c}$, degree $c+1$ if $T(0)$ is constant. That is an upper bound of degree $c+1$, not a reduction of a given degree $K$ by a factor $b<1$.

\textbf{COMPUTATIONAL} check, toy degrees only: for $K=6$, $c=2$, $C=1$, the inequality $n^{6}\ge n^{2}+(n-1)^{6}$ holds at every integer $n$ from $1$ through $12$ (equality at $n=1$). This does not say that 3-SAT has exponent $6$.

\subsection*{What a genuine contraction would require}
A charged map of the shape $a+bK$ with $0<b<1$ needs one of the following, neither of which is supplied by the self-reduction above.
\begin{enumerate}
\item[(i)] \textbf{Size shrinkage.} Instances whose bit length is at most $m^{b}$ with $b<1$. Then one call costs at most $(m^{b})^{K}=m^{bK}$, and $m^{a}$ such calls cost at most $O(m^{a+bK})$. This is the compression regime of \S\ref{sec:lit}. It is not achieved by deleting one variable.
\item[(ii)] \textbf{A non-size-reducing transformation} whose overhead has exponent $a$ and whose inherited exponent is $bK$ with $b<1$, with every cost charged, uniformly in the unknown finite $K$. This is the missing algorithm transformation. \textbf{OPEN.} No such transformation was found.
\end{enumerate}

\section{Literature}
\label{sec:lit}

\subsection*{Instance compression: conditional, and not a ban on every exponent}
\textbf{KNOWN THEOREM}, consequence \textbf{CONDITIONAL}.

Fortnow--Santhanam (ECCC TR07-096, abstract as opened): a language $L$ in NP is instance compressible if there is a polynomial-time function $f$ and a set $A$ such that, for each instance $x$, $f(x)$ has size polynomial in the \emph{witness size} of $x$ and $f$ reduces $L$ to $A$. They prove SAT is not instance compressible unless NP is contained in coNP/poly and the polynomial hierarchy collapses. The same abstract records the corollary that SAT is not polynomially kernelizable unless NP is contained in coNP/poly.

What that forbids, under the collapse hypothesis: a polynomial-time reduction whose output length is polynomial in the witness size (for SAT, polynomial in the number of variables), independent of the full input length. What it does not forbid: a reduction whose output is polynomial in the full input length (the identity is one); a kernel superpolynomial in the witness size; an unconditional prohibition of every finite decision exponent $K$; a non-shrinking transformation of the kind in (ii).

Dell--van Melkebeek (STOC 2010, Theorem~1 as opened): let $d\ge 3$ be an integer and $\varepsilon>0$. If $\mathrm{coNP}\not\subseteq\mathrm{NP}/\mathrm{poly}$, there is no oracle-communication protocol of cost $O(n^{d-\varepsilon})$ deciding satisfiability of an $n$-variable $d$-CNF, even if the first player is conondeterministic. The cost is the number of bits sent by the polynomial-time player. The same paper states a trivial protocol of cost $O(n^{d})$, tight at $\varepsilon=0$, and that a protocol of cost $O(n^{d-\varepsilon})$ would put coNP in NP/poly and collapse the polynomial hierarchy to its third level. Corollary~1, same paper: under $\mathrm{coNP}\not\subseteq\mathrm{NP}/\mathrm{poly}$ there is no polynomial-time reduction from $d$-SAT making at most $O(n^{b})$ queries, each of bit length $O(n^{c})$, whenever $b+c<d$.

What that forbids, conditionally: bit cost $O(n^{d-\varepsilon})$ for every $\varepsilon>0$, including polynomial-time sparsification, kernelization, or lossy compression down to that cost. What it does not forbid: cost $O(n^{d})$; a decision algorithm that never shrinks the instance; an unconditional time lower bound of the form $m^{K}$ against arbitrary polynomial-space 3-SAT. In particular, $O(n^{d-\varepsilon})$ is still polynomial in the number of variables. It is not, by itself, an instance of size $n^{b}$ with $b<1$ in the variable count. Relative to a dense formula, the paper notes that deleting duplicate clauses leaves $m=O(n^{d})$ clauses and that the characteristic vector has bit length $O(n^{d})$. An output of bit length $O(n^{d-\varepsilon})$ would be $O(N^{1-\varepsilon/d})$ if one set $N=\Theta(n^{d})$, and $1-\varepsilon/d<1$. That would be regime~(i) \emph{if the reduction existed}. The theorem says it does not, unless coNP $\subseteq$ NP/poly. The prohibition is of that compression, not of every conceivable non-shrinking map, and it is \textbf{CONDITIONAL}.

\subsection*{Sparsification is exponential, not a polynomial contraction}
\textbf{KNOWN THEOREM.} Impagliazzo--Paturi--Zane (author PDF, Corollary~1 of the Sparsification Lemma): for every $\varepsilon>0$ and every positive integer $k$, any $k$-CNF on $n$ variables is the disjunction of $t\le 2^{\varepsilon n}$ many $k$-CNFs, each with at most $C n$ clauses, and the disjunction is computable in time $\mathrm{poly}(n)\,2^{\varepsilon n}$. The lemma targets algorithms of the shape $2^{c n}$. The overhead is exponential in $n$. It does not send an unknown polynomial exponent $K$ to $a+bK$ with $0<b<1$, and it is not used here as such a map.

\subsection*{Alternation trading and time-space bounds restrict space}
\textbf{KNOWN THEOREM}, space-restricted. Not a map on arbitrary polynomial-time, or arbitrary polynomial-space, 3-SAT algorithms.

Williams, ECCC TR07-036 (abstract as opened): the same time-space limitation proved there for certain modular counting problems holds for SAT. The abstract states that the bound applies to time $n^{c}$ and space $n^{o(1)}$ on RAMs for $c>1$ satisfying $c^{3}-c^{2}-2c+1<0$, and that $c<1.801$ suffices. Those decimals are the abstract's comparison; they were not recomputed here. Williams' open-problems column (Gasarch column, page as opened) states the state of the art as: there is no algorithm for SAT running in time $n^{2\cos(\pi/7)-o(1)}$ and space $n^{o(1)}$ simultaneously, and writes $2\cos(\pi/7)<1.802$. Again that comparison is quoted, not recomputed. The column also records that alternation-trading proofs do not exceed this exponent (Buss--Williams). The argument simulates space-bounded RAMs. It does not take an arbitrary algorithm of time $m^{K}$ and return one of time $m^{a+bK}$.

Fortnow--Lipton--van Melkebeek--Viglas, JACM 2005 (abstract as opened): for any constant $c$ less than the golden ratio there is a positive constant $d$ such that no deterministic random-access Turing machine solves satisfiability in time $n^{c}$ and space $n^{d}$, where $d$ approaches $1$ when $c$ does. On conondeterministic machines the same holds for any constant $c$ less than $\sqrt{2}$. ECCC TR00-028 (abstract as opened) states the deterministic subpolynomial-space bound as time $n^{1.618}$ and space $n^{o(1)}$; that decimal is the abstract's figure, not a recomputation. Both results are time-space tradeoffs. They do not forbid every pair $(\mathrm{time}\ n^{K},\ \mathrm{space}\ n^{O(1)})$, and they are not a contraction of an unknown polynomial exponent in the unrestricted model.

\subsection*{Nearest self-improvement theorem, and why it is not the map}
Williams, STOC 2024, ``Self-Improvement for Circuit-Analysis Problems'' (author PDF as opened) proves a bootstrapping theorem for Circuit-SAT, \#Circuit-SAT, and quantified Circuit-SAT in an exponential / large-circuit regime: a faster algorithm for very large circuits implies a speed-up for smaller circuits, with consequences for the exponential-time hypothesis. Section~5 of that paper asks, as an open problem, whether self-improvement could go all the way down to $P=NP$, and says the current argument has bottlenecks below subexponential time. That paper does not state a uniform charged map sending every exact polynomial-time 3-SAT algorithm of exponent $K$ to one of exponent $a+bK$ with $0<b<1$. It is not cited as such a map.

No paper opened or returned by the searches states that missing map. The statement ``there exist $a$ and $b$ with $0<b<1$ such that every exact polynomial 3-SAT algorithm of exponent $K$ yields one of exponent at most $a+bK$, all costs charged'' is \textbf{OPEN}. It is not asserted.

\section{Finite check}
\label{sec:finite}
\textbf{COMPUTATIONAL}, and only for the integers printed below. Finite $n$ does not witness an asymptotic exponent.

Take the trivial exact brute-force decision cost $B(n)=2^{n}$ (one unit per assignment; a polynomial factor in $n$ would not change the conclusion below). One self-reduction step was applied in exact integer arithmetic for $n=1,\ldots,8$:
\begin{center}
\begin{tabular}{rrrr}
$n$ & $B(n)=2^{n}$ & $2\cdot B(n-1)$ & $n\cdot B(n-1)$ \\
\hline
1 & 2 & 2 & 1 \\
2 & 4 & 4 & 4 \\
3 & 8 & 8 & 12 \\
4 & 16 & 16 & 32 \\
5 & 32 & 32 & 80 \\
6 & 64 & 64 & 192 \\
7 & 128 & 128 & 448 \\
8 & 256 & 256 & 1024 \\
\end{tabular}
\end{center}
The two-call step equals $B(n)$ at each of these $n$, because $2\cdot 2^{n-1}=2^{n}$. The $n$-call step equals $n\cdot 2^{n-1}$, which at $n=8$ is $1024$, still exponential in $n$, and larger than $B(n)$ for every $n$ from $3$ through $8$. Neither column is a polynomial in $n$, so neither produces a smaller polynomial degree. The starting algorithm is not polynomial. That is the whole finite check; no large search was run.

\section*{Obstruction}
To get a contradiction one still needs a charged uniform map with $0<b<1$ whose fixed point lies in a forbidden region of the same model. Standard self-reduction has $b=1$: $n$ calls on instances of length at most $m$ send $K$ to $K+1$, and iteration of that map diverges, which is \textbf{PROVED/EXACT} as algebra and \textbf{FALSIFIED/RETIRED} as a contraction. A single call on $n-1$ variables does not force a smaller degree. A genuine contraction needs either instances of size $m^{b}$ with $b<1$, or a non-size-reducing transformation whose inherited exponent is $bK$ with $b<1$. Shrinkage strong enough to force the first alternative, at the scale Fortnow--Santhanam and Dell--van Melkebeek actually prove, is ruled out only \textbf{CONDITIONAL} on the polynomial hierarchy not collapsing, and even that rules out compression rather than every non-shrinking transformation. Impagliazzo--Paturi--Zane sparsification moves exponential $2^{\varepsilon n}$ algorithms to sparse instances; it does not yield $a,b$. Williams and Fortnow--Lipton--van Melkebeek--Viglas do not close the gap because they restrict space. No uniform charged map $K\mapsto a+bK$ with $0<b<1$ for exact 3-SAT was found. Status of the contraction: \textbf{OPEN}. Status of P versus NP: \textbf{OPEN}.

% END_OF_VERIFIED_REPORT::GROK::CONTRACT1
\end{document}
